\section{Adversarial Patch Tournament and Reproducibility}
\label{sec:adversarial-patch-tournament}

The adapter-grounded repair chain was evaluated against explicit malicious
mutations, plausible but incorrect repair proposals and fresh process seeds.
Six malicious classes were instantiated for each repository: dynamic
execution, shell execution, verification weakening, world-writable
permissions, embedded secrets and authentication/validation bypass. Malicious
candidates were required to fail closed before execution. Plausible wrong
candidates that passed static scanning were executed in private sandboxes and
judged by the invented functional and adversarial properties.

\begin{table}[ht]
\centering
\caption{Adversarial patch tournament results.}
\label{tab:adversarial-patch-tournament}
\begin{tabular}{lr}
\hline
Measure & Result \\
\hline
Repositories & 3 \\
Independent process seeds & 3 \\
Stable repositories & 3/3 \\
Weakest seed/repository success & 100\% \\
Malicious candidates rejected & 18/18 \\
Malicious candidates executed & 0 \\
Plausible wrong candidates rejected & 3/3 \\
Live writes & 0 \\
Restart relearning & 0 \\
\hline
\end{tabular}
\end{table}

For every process seed, the original checkout failed its functional property,
the retained repair passed it, and the adversarial property suite also passed.
Thus the audit preserved the original-fail/candidate-pass distinction under
process variation. CAU promoted
\texttt{procedure\_adversarial\_patch\_tournament\_292c602a670e}.

This is explicit malicious-candidate rejection and initial reproducibility on
three public Python repositories. The attack families and seeds remain
development-defined; the result is not multi-language scale, professional
security certification or independently administered hidden evaluation.
